\documentclass[11pt, letterpaper]{article}

% --- Idioma y codificación ---
\usepackage[spanish]{babel}
\usepackage[utf8]{inputenc}
\usepackage[T1]{fontenc}
\usepackage{lmodern}

% --- Matemáticas ---
\usepackage{amsmath, amssymb, amsthm}

% --- Formato ---
\usepackage[margin=2.5cm]{geometry}
\usepackage{parskip}
\usepackage{booktabs}
\usepackage{hyperref}
\hypersetup{colorlinks=true, linkcolor=blue!70!black, urlcolor=blue!70!black}

% --- Entornos de teoremas ---
\theoremstyle{definition}
\newtheorem{definition}{Definición}[section]
\theoremstyle{plain}
\newtheorem{theorem}{Teorema}[section]
\newtheorem{lemma}[theorem]{Lema}
\newtheorem{corollary}[theorem]{Corolario}
\theoremstyle{remark}
\newtheorem{remark}{Observación}[section]

\title{%
  T02 --- Demostración formal:\\[4pt]
  Cota de altura de árboles AVL
}
\author{Curso Linux--C--Rust --- B15 Estructuras de datos en C}
\date{}

\begin{document}
\maketitle
\tableofcontents

% ===================================================================
\section{Definiciones}
% ===================================================================

\begin{definition}[Árbol AVL]
Un árbol binario de búsqueda es AVL si para todo nodo $x$:
\[
  |\operatorname{h}(\operatorname{izq}(x)) - \operatorname{h}(\operatorname{der}(x))| \leq 1
\]
donde $\operatorname{h}(\text{NULL}) = -1$.
\end{definition}

\begin{definition}[Árbol de Fibonacci]
El árbol de Fibonacci de altura $h$, denotado $T_h$, es el AVL con el
\emph{mínimo} número de nodos para altura $h$:
\begin{itemize}
  \item $T_0$: un solo nodo (1 nodo, altura 0).
  \item $T_1$: raíz con un hijo izquierdo (2 nodos, altura 1).
  \item $T_h$ ($h \geq 2$): raíz con subárbol izquierdo $T_{h-1}$
    y subárbol derecho $T_{h-2}$.
\end{itemize}
\end{definition}

\begin{definition}
Sea $N(h) = |T_h|$ el número mínimo de nodos en un AVL de altura $h$.
\end{definition}

% ===================================================================
\section{Demostración 1: recurrencia de \texorpdfstring{$N(h)$}{N(h)}}
% ===================================================================

\begin{theorem}\label{thm:recurrence}
$N(h) = N(h-1) + N(h-2) + 1$ para $h \geq 2$, con $N(0) = 1$ y
$N(1) = 2$.
\end{theorem}

\begin{proof}
Sea $T$ un AVL de altura $h$ con el mínimo de nodos. Para tener
altura $h$, al menos un subárbol tiene altura $h-1$. La condición AVL
exige que el otro tenga altura $\geq h-2$. Para minimizar nodos,
elegimos exactamente $h-2$.

El subárbol de altura $h-1$ con mínimo de nodos tiene $N(h-1)$ nodos;
el de altura $h-2$ tiene $N(h-2)$. Sumando la raíz:
\[
  N(h) = N(h-1) + N(h-2) + 1
\]

Casos base: $N(0) = 1$ (un nodo, altura 0), $N(1) = 2$ (raíz $+$ un
hijo, altura 1).
\end{proof}

\begin{table}[h]
\centering
\begin{tabular}{ccl}
\toprule
$h$ & $N(h)$ & Cálculo \\
\midrule
0 & 1 & base \\
1 & 2 & base \\
2 & 4 & $2 + 1 + 1$ \\
3 & 7 & $4 + 2 + 1$ \\
4 & 12 & $7 + 4 + 1$ \\
5 & 20 & $12 + 7 + 1$ \\
6 & 33 & $20 + 12 + 1$ \\
\bottomrule
\end{tabular}
\end{table}

% ===================================================================
\section{Demostración 2: \texorpdfstring{$N(h) \geq \varphi^h$}{N(h) >= phi\^h}}
% ===================================================================

\begin{lemma}\label{lem:phi}
Sea $\varphi = \frac{1+\sqrt{5}}{2} \approx 1.618$. Entonces
$N(h) \geq \varphi^h$ para todo $h \geq 0$.
\end{lemma}

\begin{proof}
Por inducción fuerte.

\textbf{Base.}
\[
  N(0) = 1 \geq 1 = \varphi^0 \quad \checkmark
  \qquad
  N(1) = 2 \geq 1.618 = \varphi^1 \quad \checkmark
\]

\textbf{Paso inductivo ($h \geq 2$).} Suponemos $N(j) \geq \varphi^j$
para todo $j \leq h-1$:
\[
  N(h) = N(h{-}1) + N(h{-}2) + 1 > N(h{-}1) + N(h{-}2)
  \geq \varphi^{h-1} + \varphi^{h-2}
  = \varphi^{h-2}(\varphi + 1)
\]

Como $\varphi$ es raíz de $x^2 = x + 1$, se cumple
$\varphi + 1 = \varphi^2$. Verificación:
\[
  \varphi^2 = \frac{(1{+}\sqrt{5})^2}{4} = \frac{6 + 2\sqrt{5}}{4}
  = \frac{3+\sqrt{5}}{2}
  = \frac{1+\sqrt{5}}{2} + 1 = \varphi + 1 \quad \checkmark
\]

Sustituyendo:
\[
  N(h) > \varphi^{h-2} \cdot \varphi^2 = \varphi^h \qedhere
\]
\end{proof}

% ===================================================================
\section{Demostración 3: cota de altura
  \texorpdfstring{$h \leq 1.44 \cdot \log_2(n+2)$}{h <= 1.44 log2(n+2)}}
% ===================================================================

\begin{theorem}\label{thm:height}
Todo árbol AVL con $n$ nodos tiene altura
$h \leq 1.4405 \cdot \log_2(n+2)$.
\end{theorem}

\begin{proof}
Sea $T$ un AVL con $n$ nodos y altura $h$:
\[
  n \geq N(h) \geq \varphi^h
\]

Tomando logaritmo en base $\varphi$ ($\varphi > 1$):
\[
  h \leq \log_\varphi n
  = \frac{\log_2 n}{\log_2 \varphi}
  \approx \frac{\log_2 n}{0.6942}
  \approx 1.4405 \cdot \log_2 n
\]

\textbf{Cota más precisa.} Usando $N(h) = F(h{+}3) - 1$
(Demostración~4) y Binet: $N(h) \geq \varphi^{h+3}/\sqrt{5} - 2$.
De $n \geq N(h)$:
\[
  n + 2 \geq \frac{\varphi^{h+3}}{\sqrt{5}}
  \quad\Longrightarrow\quad
  h + 3 \leq \log_\varphi\!\left((n{+}2)\sqrt{5}\right)
\]
\[
  h \leq \log_\varphi(n{+}2) + \underbrace{\log_\varphi(\sqrt{5}) - 3}_{\approx\, -1.328}
  < 1.4405 \cdot \log_2(n{+}2) \qedhere
\]
\end{proof}

% ===================================================================
\section{Demostración 4: \texorpdfstring{$N(h) = F(h+3) - 1$}{N(h) = F(h+3) - 1}}
% ===================================================================

\begin{theorem}\label{thm:fib}
$N(h) = F(h+3) - 1$ para todo $h \geq 0$, donde
$F(k)$ es el $k$-ésimo número de Fibonacci con $F(1)=F(2)=1$.
\end{theorem}

\begin{proof}
Por inducción fuerte.

\textbf{Base.}
\[
  N(0) = 1 = F(3) - 1 = 2 - 1 \quad \checkmark
  \qquad
  N(1) = 2 = F(4) - 1 = 3 - 1 \quad \checkmark
\]

\textbf{Paso inductivo ($h \geq 2$).} Suponemos
$N(j) = F(j{+}3) - 1$ para todo $j \leq h{-}1$:
\begin{align*}
  N(h) &= N(h{-}1) + N(h{-}2) + 1 \\
       &= [F(h{+}2) - 1] + [F(h{+}1) - 1] + 1 \\
       &= F(h{+}2) + F(h{+}1) - 1 \\
       &= F(h{+}3) - 1 \qedhere
\end{align*}
\end{proof}

\begin{table}[h]
\centering
\begin{tabular}{ccc}
\toprule
$h$ & $F(h{+}3)$ & $N(h) = F(h{+}3) - 1$ \\
\midrule
0 & 2 & 1 \\
1 & 3 & 2 \\
2 & 5 & 4 \\
3 & 8 & 7 \\
4 & 13 & 12 \\
5 & 21 & 20 \\
6 & 34 & 33 \\
\bottomrule
\end{tabular}
\end{table}

% ===================================================================
\section{Consecuencia: operaciones AVL son \texorpdfstring{$O(\log n)$}{O(log n)}}
% ===================================================================

\begin{corollary}
Búsqueda, inserción y eliminación en un AVL con $n$ nodos tienen
complejidad $O(\log n)$ en el peor caso.
\end{corollary}

\begin{proof}
Cada operación recorre un camino de longitud $\leq h$, con trabajo
$O(1)$ por nivel. Por el Teorema~\ref{thm:height}:
\[
  h + 1 \leq 1.44 \cdot \log_2(n{+}2) = O(\log n)
\]
Las rotaciones de rebalanceo son $O(h)$ en total, cada una en $O(1)$.
\end{proof}

% ===================================================================
\section{Resumen de resultados}
% ===================================================================

\begin{table}[h]
\centering
\begin{tabular}{ll}
\toprule
\textbf{Resultado} & \textbf{Técnica} \\
\midrule
$N(h) = N(h{-}1) + N(h{-}2) + 1$
  & Minimalidad: subárboles con diferencia 1 \\
$N(h) \geq \varphi^h$
  & Inducción fuerte + $\varphi^2 = \varphi + 1$ \\
$h \leq 1.44 \cdot \log_2(n{+}2)$
  & Logaritmo de $n \geq \varphi^h$ + cambio de base \\
$N(h) = F(h{+}3) - 1$
  & Inducción fuerte + recurrencia Fibonacci \\
Operaciones AVL: $O(\log n)$
  & Corolario de la cota de altura \\
\bottomrule
\end{tabular}
\end{table}

\end{document}
